\documentclass[11pt,a4paper]{article}
% Packages for equations, images and citations
\usepackage{amsmath}
\usepackage{graphicx}
\usepackage{natbib}
\usepackage[margin=1in]{geometry}
\title{A Simple Performance Evaluation of a Classifier}
\author{Belan Bapir\\Department of Computer Engineering\\Tishk International University}
\date{December 9, 2026}
\begin{document}
\maketitle
\begin{abstract}
This teaching example explains how to calculate classification accuracy using hypothetical data. It also demonstrates the structure of a scientific paper and automatic cross-references.
\end{abstract}
\section{Introduction}
\label{sec:intro}
% Textual and parenthetical citations
\citet{Sokolova2009} discuss performance measures for classification.
Classification assigns samples to categories. Section~\ref{sec:method} explains the evaluation method, and Section~\ref{sec:results} presents an illustrative result.
\section{Methodology}
\label{sec:method}
\subsection{Evaluation Data}
Assume a classifier is tested on 100 samples, of which 85 are classified correctly.
\subsection{Performance Metric}
Accuracy is one of several measures for evaluating classifiers \citep{Fawcett2006}.
Accuracy is defined as
\begin{equation}
\label{eq:accuracy}
A=\frac{N_{\mathrm{correct}}}{N_{\mathrm{total}}}\times 100\% .
\end{equation}
\section{Results}
\label{sec:results}
Using Equation~\eqref{eq:accuracy}, the illustrative accuracy is
\begin{equation}
\label{eq:result}
A=\frac{85}{100}\times 100\%=85\%.
\end{equation}
% Figure: caption first, then label for automatic referencing.
Figure~\ref{fig:accuracy} illustrates the classification results.
\begin{figure}[htbp]
  \centering
  \includegraphics[width=0.45\textwidth]{accuracy.png}
  \caption{Illustrative classification results.}
  \label{fig:accuracy}
\end{figure}

% Table: l = left, c = center, r = right; & separates columns.
Table~\ref{tab:results} summarizes the same hypothetical data.
\begin{table}[htbp]
  \centering
  \caption{Classification performance.}
  \label{tab:results}
  \begin{tabular}{|l|c|r|}
    \hline
    Outcome & Samples & Percentage \\
    \hline
    Correct & 85 & 85\% \\
    Incorrect & 15 & 15\% \\
    Total & 100 & 100\% \\
    \hline
  \end{tabular}
\end{table}

\section{Conclusion}
Equation~\eqref{eq:result} shows that the classifier correctly identifies 85 out of 100 samples.
\section*{Acknowledgements}
Prepared as a teaching example for the TIU national LaTeX workshop. The numerical values are hypothetical, not research findings.
% Generate the reference list from references.bib.
\bibliographystyle{plainnat}
\bibliography{references}
\end{document}
